% Binary and symbolic results in III-A--III-C: 2026-09-14 artifact run.
% Source revision: 6e61076262dd06a711778489a0dea55af890ff40.
% Evidence: aisec-invariants-mlir/docs/research/satml-section3-results-20260914.agents.md.
% III-E incorporates the 2026-09-14 Overleaf contribution; its static comparison
% and regression claims were not rerun in the artifact batch above.
% This file uses only packages already present in the paper.
\section{Case Studies and Analysis Boundaries}
\label{sec:evaluation}

These cases show how lowering changes observable behaviour, why correctness and leakage require separate checks, and what each analysis can establish.

We compare each kernel on two chosen secret inputs while keeping its public inputs fixed.
Callgrind records the total numbers of executed instructions (Ir), conditional branches (Bc), and memory write references (Dw) around the kernel call, excluding input loading and program setup.
These execution counts do not measure elapsed time.
In separate runs, Memcheck tracks the secret input and reports when it influences a conditional jump or move (\texttt{cf}) or a memory address (\texttt{addr}).
Thus counts and secret-dependent conditions or addresses are separate observations.
Each count comparison uses five paired runs under the measurement procedure described in Section~IV.

We use MLIR and Clang 18.1.3, Valgrind 3.22.0, and an AMD Ryzen AI 9 HX 370 running Linux x86-64.
Core kernels follow the baseline \texttt{linalg}--\texttt{scf}--\texttt{cf}--LLVM lowering (P0), compiled with \texttt{-mavx2 -mno-avx512f} at \texttt{-O0}, \texttt{-O2}, and \texttt{-O3}.

\subsection{Select and Backend Lowering}
\label{sec:evaluation:select}

The \texttt{mask\_select} kernel selects between two public arrays using a secret byte mask.
It loads both candidate values and applies \texttt{arith.select} in a fixed-size traversal, without branching on the mask.
The lowered LLVM IR retains \texttt{select}; the secret-dependent branch appears during backend compilation at \texttt{-O0} (Figure~\ref{fig:select-lowering}).

\begin{figure}[!ht]
\centering
\begin{minipage}{0.97\columnwidth}
\footnotesize
\setlength{\topsep}{2pt}
\textbf{Source MLIR and lowered LLVM IR}
\begin{verbatim}
%c = arith.cmpi ne, %m, %zero : i8
%s = arith.select %c, %av, %bv : f32
...
%39 = select i1 %38, float %34, float %37
\end{verbatim}
\textbf{O0: the mask controls a jump}
\begin{verbatim}
3550: cmp $0x0,%al
3558: jne 356a <mask_select+0x9a>
\end{verbatim}
\textbf{O2/O3: the vector path uses a blend}
\begin{verbatim}
2df2: vpcmpeqd %ymm0,%ymm1,%ymm1
2df6: vblendvps %ymm1,(%rdx,%rax,4),
                %ymm5,%ymm1
\end{verbatim}
\end{minipage}
\caption{Measured select lowering. Surrounding loads and public loop/alias checks are omitted; the blend instruction is wrapped. LLVM IR retains selection until backend compilation.}
\label{fig:select-lowering}
\end{figure}
% Coauthor note kept from the Overleaf version:
% \EK{I believe we could remove code freely to save space - we have 13 pages already}

% Author question addressed by the measurement paragraph above:
% \SP{We should give some background on what is being measured and how}
At \texttt{-O0}, the all-one mask executes 8,192 fewer instructions and 4,096 fewer write references than the all-zero mask.
Both executions encounter the same number of conditional branches, but Memcheck identifies a branch whose condition depends on the mask.
Equal branch counts therefore do not imply equal branch outcomes.

% Author question addressed by specifying the kernel call below:
% \SP{What is a marked region?}
The exercised vector path uses \texttt{vblendvps} at \texttt{-O2} and \texttt{-O3}.
Memcheck reports no secret-dependent conditional jump, conditional move, or memory address during the kernel call, and the two masks produce identical instruction, branch, and write-reference counts.
This result applies to the tested executions; it is not a proof of constant-time behaviour.

% \SP{what is a conditional reduction?}
Table~\ref{tab:section3-core} compares the select with a fixed computation, a branch on a secret sum, a load at a secret index, and a loop with a secret extent.
The \texttt{cond\_reduce} kernel sums its secret array, then branches on whether the sum is positive to write either 1 or 2.

\begin{table}[t]
\caption{P0 results for the chosen inputs. Memcheck reports \texttt{cf} or \texttt{addr}; count difference denotes a stable difference in Ir, Bc, or Dw. A dash means no finding or count difference detected, respectively.}
\label{tab:section3-core}
\centering
\footnotesize
\setlength{\tabcolsep}{3pt}
\begin{tabular}{@{}lcccccc@{}}
\hline
& \multicolumn{3}{c}{Memcheck finding} & \multicolumn{3}{c}{Count difference} \\
\cline{2-4}\cline{5-7}
Kernel & O0 & O2 & O3 & O0 & O2 & O3 \\
\hline
\texttt{matvec} & --- & --- & --- & --- & --- & --- \\
\texttt{cond\_reduce} & \texttt{cf} & --- & --- & Yes & --- & --- \\
\texttt{mask\_select} & \texttt{cf} & --- & --- & Yes & --- & --- \\
\texttt{idx\_gather} & \texttt{addr} & \texttt{addr} & \texttt{addr} & --- & --- & --- \\
\texttt{dynshape} & \texttt{cf} & \texttt{cf} & \texttt{cf} & Yes & Yes & Yes \\
\hline
\end{tabular}
\end{table}

\subsection{Shape and Sparsity}
\label{sec:evaluation:structure}

% \SP{stored-value counts?}
Confidential structure can influence execution even when tensor dimensions and the number of explicitly stored entries are fixed.
Our sparse example represents a vector of length 256 by eight stored values \(1,\ldots,8\), their coordinates, and a position array \([0,8]\).
% \SP{What is the difference between a position and a coordinate?}
% \SP{what is a public position array?}
% \SP{What is a coordinate, still unclear}
The positions are offsets into the stored-entry arrays: \([0,8]\) selects entries 0 through 7.
Each coordinate gives the index in the length-256 vector where its associated value belongs.
The vector length, values, and positions are public, meaning the observer may know them, and are identical in both inputs.
Only the coordinate array changes: the first input uses \([0,1,2,3,4,5,6,7]\), while the second uses \([3,29,61,97,130,168,201,240]\).

\paragraph{From representation to address}
The source first assembles values, positions, and coordinates into a compressed sparse tensor, then converts it to a dense tensor.
We apply MLIR's sparsifier with its runtime-library mode disabled, followed by LLVM conversion and a Clang \texttt{-O0} backend.
The code zeroes a fixed-size output, then iterates over the eight stored entries.
The public position array determines the loop bounds; each secret coordinate determines a store address (Figure~\ref{fig:sparse-address}).

\begin{figure}[!ht]
\centering
\begin{minipage}{0.97\columnwidth}
\footnotesize
\setlength{\topsep}{2pt}
\textbf{Dense materialization, summarized as pseudocode}
\begin{verbatim}
out[0:256] = 0
for i in pos[0]:pos[1]:     # public range
    c = crd[i]             # secret coordinate
    out[c] = vals[i]       # secret address
\end{verbatim}
\textbf{Generated address and store, with SSA names shortened}
\begin{verbatim}
%c = llvm.load %crd_ptr : !llvm.ptr -> i64
%p = llvm.getelementptr %out[%c]
     : (!llvm.ptr, i64) -> !llvm.ptr, f32
llvm.store %v, %p : f32, !llvm.ptr
\end{verbatim}
\textbf{Machine instruction reported by Memcheck}
\begin{verbatim}
1976: movss %xmm0,(%rcx,%rdx,4)
\end{verbatim}
\end{minipage}
\caption{The sparse-to-dense lowering uses a secret coordinate as a store index. The public number of iterations does not determine the sequence of addresses.}
\label{fig:sparse-address}
\end{figure}

\paragraph{Equal counts, different addresses}
The scatter loop performs eight stores for either coordinate pattern, but their destination addresses differ.
The sparse binary has equal instruction, branch, and write-reference counts in all five paired runs, while Memcheck reports a store address that depends on the secret coordinate.
The \texttt{idx\_gather} example shows the same distinction at all three optimisation levels.
Equal counts do not rule out secret-dependent addresses, though these experiments do not measure cache-line visibility or physical timing.

The separately authored dense reference visits all \(256\times8\) position-entry pairs and uses loop-derived load and store addresses.
No secret-dependent address is detected in that build, but its \texttt{-O0} binary contains a secret-dependent \texttt{je}, confirmed at the instruction identified by Memcheck.
Its execution counts also remain equal across the two coordinate inputs.
The comparison distinguishes their address channels; the reference already has a control channel.
As these are separate source programs without a complete output comparison, the experiment does not establish functional correctness of the sparse lowering.

\paragraph{Dynamic extent and bufferization}
The \texttt{dynshape} kernel makes the allocation size and loop bound depend on a secret extent \(k\).
% \SP{What is a backend level, do you mean opt level?}
For \(k=1\) versus \(k=4096\), Memcheck reports a secret-dependent condition at all three P0 optimisation levels.
The instruction difference falls from 73,872 at \texttt{-O0} to 5,117 at \texttt{-O2}/\texttt{-O3}; the write-reference difference falls from 20,493 to zero as the write-only buffer is eliminated.
This experiment starts from \texttt{memref} operations.
Attributing a channel to one-shot bufferization requires comparison with the corresponding tensor program lowered through that pass (P4).

\subsection{Two Distinct Verification Obligations}
\label{sec:evaluation:polygeist}

A transformation can compute the wrong result without adding leakage, or preserve the result while adding an observation.
We check four hand-transcribed Polygeist transformation specifications: two faithful examples and two deliberately modified versions.
The equivalence check compares returned values and modeled final memory under the checker's refinement relation.
The leakage-preservation check asks whether source-indistinguishable executions remain target-indistinguishable:
% \SP{What do these symbols mean?}
\[
  L_S(x)=L_S(x')\ \Longrightarrow\ L_T(x)=L_T(x').
\]
Here \(L_S\) and \(L_T\) are the source and target observation traces, and \(x,x'\) agree on public inputs.
The source may already leak; the target must not distinguish inputs the source's observations leave indistinguishable.
The two obligations are logically distinct but share the SMT encoding infrastructure.
Table~\ref{tab:gate-controls} uses the supported observation model with the \texttt{--unroll 8} truncation limit.
The faulty specifications are deliberate modifications, not discovered upstream compiler bugs.


Propagating an unchanged loop-carried value passes both checks.
The modified specification substitutes its initial value after it changes: the returned accumulator exposes the error, while the leakage check passes.
The faithful loop restructuring also passes both checks.
Moving its body before the first condition check fails leakage preservation but passes equivalence for this example's returned flag.
The flag is false within the modeled executions; this does not establish that arbitrary loop bodies can be moved safely.

Verification succeeds only when both checks pass; each rejected specification returned a counterexample.
All four results concern the model truncated with \texttt{--unroll 8}; the saved loop counterexample is sensitive to that truncation.
For CFG loops, the limit counts visits per block rather than restricting input bounds to eight completed iterations.
Unspecified computations, called holes, are symbolic functions shared between source and target, with effects and observations limited by the model.
The results cover these specifications, not production passes, unbounded loops, or downstream executables.

\subsection{Analysis Boundaries}
\label{sec:evaluation:boundaries}

Table~\ref{tab:analysis-boundaries} compares the methods' scopes: one IR stage, symbolic inputs under a truncated model, or an executable on chosen inputs.


The select has no secret branch in the inspected IR but has one in the \texttt{-O0} executable: the compilation stage matters.
Gather and sparse materialization use secret-dependent addresses despite equal counts: the observation matters.
No method supplies universal ground truth for the others.

Unsupported encodings, timeouts, and solver unknown responses yield \texttt{UNKNOWN}; none occurred in the four verification examples.
A bounded result is not automatically \texttt{UNKNOWN}.
IR analysis reports some unsupported operations, but missing dependency propagation rules can cause missed findings.
Without validated rule coverage, an IR result with no finding does not establish obliviousness.
The comparison in Section~\ref{sec:evaluation:static} concerns five kernels under the stated lowerings; agreement on that corpus does not establish complete rule coverage or general analyser accuracy.

\subsection{Static prediction against measured ground truth}
\label{sec:evaluation:static}
% Imported comparison: verify the exact static-analysis revision, commands,
% per-pipeline outputs, regression test, and runtime before submission.

Measurement (§IV) speaks about the binary that ran, at one configuration and for two input classes.
Symbolic verification speaks about inputs within its supported execution and observation model.
A third tier sits between them: a forward taint analysis over the MLIR as given, needing no flattening, solver or per-operation semantics, and reasoning about the IR rather than the downstream binary.
An all-inputs conclusion requires sound propagation rules for the supported IR.
This subsection defines it and asks whether its verdicts agree with what the binary did.

\paragraph{What the analysis is} A function argument marked secret taints its value, and taint propagates forward: any operation with a tainted operand taints all of its results, a store taints the buffer it writes, a yield taints both the parent operation's result and the loop argument it re-binds, and a branch's successor operands taint the block arguments they reach.
These edges run in both directions around a loop, so the taint sets are iterated to a fixpoint rather than computed in one pass.
A tainted value that reaches a loop bound, a load or store index, or a branch condition is a violation on the control or address obligation.
Unsupported calls, buffer aliases, and opaque regions must produce an unknown result when their effects cannot be modeled.
The analysis uses forward taint propagation, as does the verifier's prefilter, but operates directly on unflattened IR.
This lets it inspect some programs that the flattener refuses; it does not establish that the two implementations have identical coverage.
That propagation is standard taint over ordinary MLIR annotations, and we claim nothing for it (discussed in related work).
What we claim is the comparison.

\paragraph{The comparison} The five kernels of §IV were run through the analysis before and after the MLIR pipelines described in Section~IV, at the same lowerings the rig measured, and the static verdict was set beside the measured one at -O0.
Table~\ref{tab:static-analysis-results} gives the result.

The comparison distinguishes five verdicts: \emph{oblivious}, meaning no finding before or after lowering; \emph{authored-and-survives}, a violation before and after; \emph{lowering-introduced}; \emph{lowering-removed}; and \emph{unknown-after}.
An unknown result after lowering must not be counted as removal.
The reported static verdicts are the same for P0--P3 and P5; P4 applies to the tensor kernels.
In Table~\ref{tab:static-analysis-results}, ``Authored'' abbreviates authored-and-survives.
Agreement compares the presence of a control or address finding, not equality of every observation.
The binary column retains the reported control and address findings; \texttt{Dw} denotes a difference in write-reference counts, rather than an additional static obligation.
The one-write difference for \texttt{cond\_reduce} accompanies its control finding and is omitted from this summary.

\begin{table}[t]
% \EK{Oh my Godness, wha happened}
\caption{Reported static and binary verdicts at \texttt{-O0}.}
\label{tab:static-analysis-results}
\centering
\footnotesize
\setlength{\tabcolsep}{3pt}
\begin{tabular}{@{}lccc@{}}
\hline
Kernel & Binary & IR before/after & Agree \\
\hline
\texttt{matvec} & Oblivious & Oblivious & Yes \\
\texttt{cond\_reduce} & \texttt{cf} & Authored (\texttt{cf}) & Yes \\
\texttt{idx\_gather} & \texttt{addr} & Authored (\texttt{addr}) & Yes \\
\texttt{dynshape} & \texttt{cf}+\texttt{Dw} & Authored (\texttt{cf}) & Yes \\
\texttt{mask\_select} & \texttt{cf} & Oblivious & \textbf{No} \\
\hline
\end{tabular}
\end{table}

The reported verdicts agree on four of five kernels; the disagreement has a compilation-stage explanation.
At the MLIR level \texttt{mask\_select} has no secret-dependent branch or address; \texttt{arith.select} on the mask is exactly what the analysis models, and a finding-free result is consistent with that source structure.
The measured branch appears during LLVM backend compilation at \texttt{-O0} (Section~\ref{sec:evaluation:select}), after the inspected IR stage; at -O2, where the selector emits a blend, the measured result has no finding and the two tiers agree again --- with the static tier having seen neither the introduction nor the removal.
The comparison illustrates the compilation-stage limit described in Section~\ref{sec:evaluation:boundaries}.
Agreement on these kernels is not a general accuracy or soundness guarantee.

\paragraph{A false verdict caught by measurement} The comparison did work beyond confirming agreement.
An earlier form of the analysis lost taint across \texttt{--convert-scf-to-cf}: it tracked \texttt{scf} conditions but not the \texttt{cf.cond\_br} they lower to, so one standard lowering step later it went blind.
A before/after run then read that blindness as lowering-removed --- the compiler appeared to have fixed \texttt{cond\_reduce}, whose \texttt{-O0} binary is reported to retain a control finding across the measured pipelines.
The blindness, not the compiler, had removed the leak.
The reported fix propagates the missing edges and adds a regression kernel, but the lesson generalizes to any differential analysis: a verdict of removed is trustworthy only if the analysis can still see the program after the step, and one that cannot must answer unknown-after --- which must never be folded into removed.

\paragraph{What the tier is for} It covers two of the four obligations, models no aliasing, is intraprocedural, and stops at the MLIR boundary.
Within the supported IR it can propagate dependencies without translating each operation into SMT.
Its conclusions still depend on sound propagation and detection of unsupported constructs; the five-kernel comparison does not validate those properties for arbitrary programs.
